
\subsubsection{2.1 LLM Calibration}

Guo et al. (2017) formalized Expected Calibration Error (ECE) and demonstrated that modern deep networks are systematically overconfident. Kadavath et al. (2022) extended calibration analysis to large language models, showing calibration improves with scale but varies substantially by domain. Zhao et al. (2023) demonstrated that generation-based ECE correlates with TruthfulQA hallucination rates across model families, providing the closest precedent to the present work; however, that work did not perform systematic cross-family factor analysis or explicit capability controls.

\subsubsection{2.2 Hallucination and Factual Reliability}

Lin et al. (2022) introduced TruthfulQA, demonstrating that larger language models are not necessarily more truthful. Yin et al. (2023) showed that overconfidence and under-refusal are measurable failure modes correlated with model scale. These works establish hallucination and miscalibration as related phenomena, but the quantitative population-level relationship under capability control has not been characterized.

\subsubsection{2.3 Adversarial Robustness}

Wang et al. (2021) introduced AdvGLUE, demonstrating substantial accuracy drops under adversarial perturbations even for high-accuracy models. Nie et al. (2020) developed ANLI via iterative adversarial filtering. These works establish adversarial robustness as a distinct evaluation dimension but do not analyze its correlation with calibration or hallucination across diverse model populations.

\subsubsection{2.4 Multi-Dimensional Trustworthiness Evaluation}

Liang et al. (2022) evaluate models across 42 scenarios and 7 metrics in HELM, noting that rankings differ substantially by metric but explicitly declining to analyze cross-metric correlations. Wang et al. (2023) provide comprehensive trustworthiness assessment for GPT-3.5 and GPT-4 in DecodingTrust, finding that trustworthiness properties are not uniformly correlated; their scope is limited to GPT models. Sun et al. (2024) evaluate six dimensions for 16 LLMs in TrustLLM, reporting partial calibration–truthfulness correlations without factor analysis or capability control.

\subsubsection{2.5 Psychometric Approaches to Model Evaluation}

Burnell et al. (2023) argue for more principled statistical approaches to LLM evaluation. Polo et al. (2024) apply Item Response Theory to benchmark data. The present work extends this methodological direction to multi-property trustworthiness structure analysis.

\subsubsection{2.6 Positioning}

No prior work computes a capability-controlled cross-property correlation matrix across calibration (ECE, Brier), hallucination (TruthfulQA\%), and adversarial robustness (AdvGLUE drop, ANLI drop) for a diverse open-weight model population, followed by factor analysis to test for a coherent latent dimension.

